%mac
\begin{dang}{TỔNG HỢP - PHÂN TÍCH LỰC VÀ ĐIỀU KIỆN CÂN BẰNG}
\begin{itemize}
\item Đầu tiên ta biểu diễn chính xác các vector lực tác dụng lên vật về phương, chiều và độ lớn (với độ lớn ta có thể ước lượng).
\item Điều kiện cân bằng của chất điểm: $\vec{F}=\vec{F}_1+\vec{F}_2+\vec{F}_3=\vec{0}$.
\end{itemize}
\newghichu{
\begin{itemize}
\item Lực căng (thường kí hiệu là $\vec{T}$) luôn hướng vào trong sợi dây.
\item Lực nén của thanh (thường kí hiệu $\vec{R}$ hoặc $\vec{N}$) luôn hướng ra ngoài thanh.
\item Phản lực (thường kí hiệu $\vec{N}$) đặt tại điểm tiếp xúc, phương vuông góc với mặt đỡ và hướng vào trong vật.
\end{itemize}
}
\paragraph{Phương pháp hình học - Áp dụng cho ba lực cân bằng}
\begin{itemize}
\item Chất điểm chịu ba lực cân bằng: $\vec{F}=\vec{F}_1+\vec{F}_2+\vec{F}_3=\vec{0}$
\item Ta tổng hợp cặp lực khéo léo sao cho: $\vec{F}_{12}=-\vec{F}_3$
\item Sử dụng các công thức hình học để tìm các đại lượng cần tìm.
\end{itemize}
\paragraph{Phương pháp đại số}
\begin{itemize}
\item Chọn hệ trục Oxy vuông góc (chọn sao cho phép chiếu đơn giản nhất)
\item Do chất điểm cân bằng nên nó sẽ cân bằng theo mọi phương. Ta chiếu các lực nên hai phương Ox và Oy. 
\item $\begin{cases}
F_x=F_{1x}+F_{2x}+...+F_{nx}=0\\
F_y=F_{1y}+F_{2y}+...+F_{ny}=0
\end{cases}$
\item Giải hệ tìm ra các đại lượng cần tìm.
\end{itemize}
\end{dang}

